\documentclass[11pt]{amsart}

\usepackage{amsmath,amssymb,amsthm,mathtools}
\usepackage[colorlinks=true,linkcolor=blue,citecolor=blue,urlcolor=blue]{hyperref}

\newtheorem{theorem}{Theorem}[section]
\newtheorem{proposition}[theorem]{Proposition}
\newtheorem{corollary}[theorem]{Corollary}
\newtheorem{lemma}[theorem]{Lemma}
\theoremstyle{definition}
\newtheorem{definition}[theorem]{Definition}
\newtheorem{remark}[theorem]{Remark}
\newtheorem{problem}[theorem]{Problem}

\newcommand{\N}{\mathbb{N}}
\newcommand{\R}{\mathbb{R}}
\newcommand{\Z}{\mathbb{Z}}

\title[Modular-control reformulation]{A modular-control reformulation of the regular induced subgraph conjecture}
\author{Arthur Freitas Ramos}
\address{Microsoft, USA}
\email{arfreita@microsoft.com}

\author{David Barros Hulak}
\address{Independent researcher}
\email{dbhulak@gmail.com}

\author{Ruy B.~J.~G. de Queiroz}
\address{Centro de Inform\'atica, Universidade Federal de Pernambuco, Brazil}
\email{ruy@cin.ufpe.br}

\subjclass[2020]{Primary 05C35; Secondary 05C07, 05C69}
\keywords{regular induced subgraph; modular degree; structural reformulation;
control set; asymptotic frontier}

\begin{document}

\begin{abstract}
Let $F(n)$ be the largest integer such that every graph on $n$ vertices contains a regular induced
subgraph on at least $F(n)$ vertices. We record a structural reformulation of the conjecture that
$F(n)$ grows superlogarithmically. For every integer base $b>1$, the conjecture is equivalent to
linear growth along the subsequence $b^k$, to the corresponding modular target, to a sharper
low-degree modular target in which the modulus need only exceed the maximum induced degree, and to
an exact one-control formulation with a singleton control vertex. On the modular side, we isolate
the one-control modular witness as the canonical asymptotic object and show that the natural
multiblock, bucketing, and cascade variants carry no extra asymptotic information: they are all
equivalent to the same target. Two transport lemmas and a paired-residue bucketing lemma with
quadratic cost in the modulus support this reduction on the finite side. The point is not a new
lower bound, but a proof that several apparently stronger modular-control packages collapse to the
same frontier.
\end{abstract}

\maketitle

\section{Introduction}

All graphs in this paper are finite and simple. For $n \in \N$, let $F(n)$ be the largest integer
such that every graph on $n$ vertices contains a regular induced subgraph on at least $F(n)$
vertices. The conjecture usually attributed to Erd\H{o}s, Fajtlowicz, and Staton asks whether
\[
  \frac{F(n)}{\log n} \longrightarrow \infty .
\]
The basic Ramsey argument gives only the logarithmic lower bound $F(n) \ge c \log n$ for some
absolute $c>0$; compare \cite{FajtlowiczMcColganReidStaton,AlonKrivelevichSudakov}. Further
progress in adjacent settings includes polynomial lower bounds for random graphs
\cite{KrivelevichSudakovWormald}, algorithmic and structural results for $k$-regular induced
subgraphs \cite{CardosoKaminskiLozin}, and modular-degree theorems in sparse graphs such as
the result of Berman, Radcliffe, Scott, Wang, and Wargo that every tree contains a large induced
subgraph with all degrees $\equiv 1 \pmod k$ \cite{BermanRadcliffeScottWangWargo}. Modular degree
conditions also appear in Scott's and Caro's work on induced subgraphs with prescribed residues
\cite{Scott,Caro}. These results show that congruence constraints arise naturally in induced
subgraph problems, but they do not identify the minimal witness package relevant to the regular
induced subgraph conjecture. We do not improve any of these bounds.

Instead, the aim of this paper is structural. We isolate a modular-control reformulation of the
conjecture in which the remaining open step becomes quite explicit. The contribution is not a new
modular-degree theorem, but a normalization result for the conjecture: once one allows
ambient-degree and external-degree congruences, there is no asymptotic gain in keeping several
control blocks, a dropped part, or a cascade of refinements. Those packages all collapse to a
single one-control modular witness target (Theorem~\ref{thm:canonical-frontier}). Concretely, the
paper proves three reductions:

\begin{itemize}
  \item the conjecture is equivalent to its restriction to any geometric subsequence $b^k$, and to
    both the usual modular target and the sharper low-degree modular target;
  \item on the exact side, the same power-scale target is equivalent to the existence of linearly
    large one-control exact witnesses, even with a singleton control set;
  \item on the modular side, the one-control, multiblock, bucketing, and cascade witness packages
    are all asymptotically equivalent.
\end{itemize}

The key message is therefore not merely that modular language is available, but that the conjecture
can be reduced to low-complexity degree data without sacrificing asymptotic strength. The finite
lemmas behind this reduction are simple:

\begin{itemize}
  \item if all vertices of a set $S$ have the same degree in $G[S \cup T]$ and the same degree into
    a disjoint control set $T$, then $G[S]$ is regular;
  \item the modular analogue holds with congruences modulo $q$;
  \item when one buckets by the pair consisting of the ambient degree modulo $q$ and the degree into
    $T$ modulo $q$, only $q^2$ residue classes appear.
\end{itemize}

These lemmas suggest a control-set route that is exponentially cheaper than tracking full
adjacency patterns. More importantly, they show that any successful modular-control argument can,
after the fact, be normalized to one control set and one modulus. This is the sense in which the
reformulation is intended to be canonical rather than cosmetic.

\medskip

The paper is organized as follows. Section~\ref{sec:powerscales} records the geometric-scale and
modular reformulations of the conjecture. Section~\ref{sec:transport} states the finite exact and
modular transport lemmas together with the paired-residue bucketing principle. In
Section~\ref{sec:frontier} we define the one-control exact and modular witnesses, explain why the
exact one-control target is already equivalent to the usual geometric formulation, and isolate the
one-control modular target as the canonical modular frontier together with its multiblock,
bucketing, and cascade reformulations; Theorem~\ref{thm:summary-package} then packages the paper as
a single equivalence theorem. The final section records a concrete open problem suggested by this
reorganization.

\subsection*{Notation}
All graphs are finite and simple. For a graph $G$ and a vertex set $S$, we write $G[S]$ for the
induced subgraph on $S$. For disjoint sets $S,T$ and $v \in S$ we write $\deg_{G[S]}(v)$ for the
degree of $v$ in $G[S]$, $\deg_{G[S\cup T]}(v)$ for its degree in $G[S \cup T]$, and
$\deg_T(v) = |N_G(v)\cap T|$ for its number of neighbors in $T$. Congruences $a \equiv b\pmod q$
are taken in $\Z$ in the usual sense; we say a quantity is constant modulo $q$ on a set $U$ if
its values on $U$ are pairwise congruent modulo $q$.



\section{Power scales and modular witnesses}\label{sec:powerscales}

\begin{definition}
Fix an integer base $b>1$.
\begin{enumerate}
  \item We say that the \emph{power-sequence target} holds at base $b$ if, for every $M \in \N$,
    there exists $K$ such that for every $k \ge K$, every graph on $b^k$ vertices contains a
    regular induced subgraph on at least $Mk$ vertices.
  \item A \emph{modular witness} in a graph $G$ is an induced subgraph $G[S]$ such that all degrees
    in $G[S]$ are congruent modulo some integer $q \ge |S|$.
  \item A \emph{low-degree modular witness} in a graph $G$ is an induced subgraph $G[S]$ such that
    all degrees in $G[S]$ are congruent modulo some integer $q$ and every degree in $G[S]$ is
    strictly smaller than $q$.
  \item We say that the \emph{power-sequence modular target} holds at base $b$ if, for every
    $M \in \N$, there exists $K$ such that for every $k \ge K$, every graph on $b^k$ vertices
    contains a modular witness of order at least $Mk$.
  \item We say that the \emph{power-sequence low-degree modular target} holds at base $b$ if, for
    every $M \in \N$, there exists $K$ such that for every $k \ge K$, every graph on $b^k$
    vertices contains a low-degree modular witness of order at least $Mk$.
\end{enumerate}
\end{definition}

The first reformulation is the standard but extremely useful observation that the asymptotic
statement may be tested on any fixed geometric subsequence.

\begin{theorem}\label{thm:power-sequence}
For every integer $b>1$, the following are equivalent.
\begin{enumerate}
  \item $F(n)/\log n \to \infty$.
  \item The power-sequence target holds at base $b$.
\end{enumerate}
\end{theorem}

\begin{proof}
Assume first that $F(n)/\log n \to \infty$. Fix $M \in \N$. Then
\[
  \frac{F(b^k)}{\log(b^k)} = \frac{F(b^k)}{k \log b} \longrightarrow \infty ,
\]
so for all sufficiently large $k$ we have $F(b^k) \ge Mk$. This is exactly the power-sequence
target.

Conversely, assume the power-sequence target at base $b$. Let $C>0$. It is enough to show that
$F(n) \ge C \log n$ for all sufficiently large $n$. Choose an integer $M > 2C \log b$. Let $K$ be
such that every graph on $b^k$ vertices contains a regular induced subgraph of order at least $Mk$
for all $k \ge K$. Now let $n \ge b^K$, and choose $k$ with $b^k \le n < b^{k+1}$.

By monotonicity of $F$,
\[
  F(n) \ge F(b^k) \ge Mk.
\]
Since $\log n < (k+1)\log b$, we obtain
\[
  C \log n < C(k+1)\log b \le \frac{M}{2}(k+1).
\]
For all sufficiently large $k$ one has $k+1 \le 2k$, and then
\[
  C \log n < Mk \le F(n).
\]
Therefore $F(n)$ eventually dominates every constant multiple of $\log n$, which is equivalent to
the conjecture.
\end{proof}

The modular route starts from the observation that congruence plus a sufficiently large modulus is
already exact regularity.

\begin{proposition}\label{prop:modular-collapse}
Let $G$ be a graph and let $S$ be a vertex set. If all degrees in $G[S]$ are congruent modulo an
integer $q \ge |S|$, then $G[S]$ is regular.

More generally, if all degrees in $G[S]$ lie in an interval of width strictly smaller than $q$ and
are pairwise congruent modulo $q$, then $G[S]$ is regular.
\end{proposition}

\begin{proof}
Every degree in $G[S]$ lies in $\{0,1,\dots,|S|-1\}$. If $q \ge |S|$, then two numbers in that set
that are congruent modulo $q$ must coincide. The interval version is the same observation: two
integers in an interval of width $<q$ that are congruent modulo $q$ are equal.
\end{proof}

This already converts the geometric regular-subgraph target into a modular one.

\begin{theorem}\label{thm:modular-reduction}
For every integer $b>1$, the following are equivalent.
\begin{enumerate}
  \item $F(n)/\log n \to \infty$.
  \item The power-sequence modular target holds at base $b$.
  \item The power-sequence low-degree modular target holds at base $b$.
\end{enumerate}
\end{theorem}

\begin{proof}
By Theorem~\ref{thm:power-sequence}, it is enough to compare both statements with the
power-sequence target. Every regular induced subgraph is trivially a modular witness: take any
modulus $q \ge |S|$. It is also a low-degree modular witness: if all degrees in $G[S]$ equal $d$,
take $q=d+1$.

Conversely, Proposition~\ref{prop:modular-collapse} turns every modular witness into an exact
regular induced subgraph of the same order. A low-degree modular witness is also exact by the
interval version of that proposition, since all induced degrees lie in the interval $[0,q)$ of
width $q$.
\end{proof}

\begin{remark}
The modular reformulation is not stronger than the original problem. Its value is organizational:
congruence conditions interact better with control sets and bucketing than exact equality does. The
low-degree version sharpens this further: the modulus need only dominate the induced degree range,
not the full cardinality of the surviving set.
\end{remark}

\section{Transport lemmas and cheap bucketing}\label{sec:transport}

The most basic control-set lemma is the following.

\begin{proposition}[exact control-set transport]\label{prop:exact-transport}
Let $S$ and $T$ be disjoint vertex sets in a graph $G$. Suppose that all vertices of $S$ have the
same degree in $G[S \cup T]$ and also the same degree into $T$. Then $G[S]$ is regular.
\end{proposition}

\begin{proof}
For every $v \in S$,
\[
  \deg_{G[S]}(v)=\deg_{G[S \cup T]}(v)-\deg_T(v).
\]
If the two terms on the right are independent of $v$, then so is the left-hand side.
\end{proof}

The modular analogue is equally simple.

\begin{proposition}[modular control-set transport]\label{prop:modular-transport}
Let $S$ and $T$ be disjoint vertex sets in a graph $G$, and let $q \ge 1$. If the degrees of all
vertices of $S$ in $G[S \cup T]$ are pairwise congruent modulo $q$, and the degrees of all
vertices of $S$ into $T$ are pairwise congruent modulo $q$, then the degrees in $G[S]$ are pairwise
congruent modulo $q$.
\end{proposition}

\begin{proof}
For $v,w \in S$ one has
\begin{align*}
  \deg_{G[S \cup T]}(v)-\deg_{G[S]}(v) &= \deg_T(v), \\
  \deg_{G[S \cup T]}(w)-\deg_{G[S]}(w) &= \deg_T(w).
\end{align*}
Subtracting the two identities and using the congruence hypotheses gives
$\deg_{G[S]}(v) \equiv \deg_{G[S]}(w) \pmod q$.
\end{proof}

The same bookkeeping works with several control blocks.

\begin{proposition}[multiblock modular transport]\label{prop:multiblock-transport}
Let $S,T_1,\dots,T_r$ be pairwise disjoint vertex sets in a graph $G$, and let $q \ge 1$. If the
degrees of all vertices of $S$ in $G[S \cup T_1 \cup \cdots \cup T_r]$ are pairwise congruent
modulo $q$, and for each $i=1,\dots,r$ the degrees of all vertices of $S$ into $T_i$ are pairwise
congruent modulo $q$, then the degrees in $G[S]$ are pairwise congruent modulo $q$.
\end{proposition}

\begin{proof}
For every $v \in S$,
\[
  \deg_{G[S]}(v)
  =
  \deg_{G[S \cup T_1 \cup \cdots \cup T_r]}(v)-\sum_{i=1}^r \deg_{T_i}(v).
\]
Each term on the right is constant modulo $q$ on $S$ by hypothesis, so the left-hand side is also
constant modulo $q$ on $S$.
\end{proof}

The key combinatorial gain is that one need only freeze two scalar quantities modulo $q$, rather
than the full adjacency pattern into the control set.

\begin{proposition}[paired-residue bucketing]\label{prop:q2-bucketing}
Let $S$ and $T$ be disjoint vertex sets in a graph $G$, let $q \ge 1$, and let $k \in \N$. If
$q^2 k \le |S|$, then there exists a subset $U \subseteq S$ with $|U| \ge k$ such that all degrees
in $G[S]$ on $U$ are congruent modulo $q$.

More precisely, one may choose $U$ so that both the ambient degree in $G[S \cup T]$ modulo $q$ and
the degree into $T$ modulo $q$ are constant on $U$.
\end{proposition}

\begin{proof}
Consider the map
\[
  \phi \colon S \to (\mathbb{Z}/q\mathbb{Z})^2,
  \qquad
  \phi(v)=\bigl(\deg_{G[S \cup T]}(v)\bmod q,\ \deg_T(v)\bmod q\bigr).
\]
Its codomain has exactly $q^2$ elements. Since $q^2 k \le |S|$, the pigeonhole principle yields a
fiber $U \subseteq S$ with $|U| \ge k$ on which $\phi$ is constant. Thus all vertices of $U$ have
the same degree in $G[S \cup T]$ modulo $q$ and the same degree into $T$ modulo $q$.

For $u,w \in U$ we have
\begin{align*}
  \deg_{G[S]}(u) &= \deg_{G[S \cup T]}(u)-\deg_T(u), \\
  \deg_{G[S]}(w) &= \deg_{G[S \cup T]}(w)-\deg_T(w),
\end{align*}
because $S$ and $T$ are disjoint. The two terms on the right are pairwise congruent modulo $q$, so
$\deg_{G[S]}(u)\equiv \deg_{G[S]}(w)\pmod q$. Hence the degrees in $G[S]$ on $U$ are all congruent
modulo $q$.
\end{proof}

The same pigeonhole argument extends to several control blocks.

\begin{proposition}[multiblock residue bucketing]\label{prop:multiblock-bucketing}
Let $S,T_1,\dots,T_r$ be pairwise disjoint vertex sets in a graph $G$, let $q \ge 1$, and let
$k \in \N$. If $q^{r+1}k \le |S|$, then there exists a subset $U \subseteq S$ with $|U| \ge k$
such that all vertices of $U$ have the same degree in $G[S \cup T_1 \cup \cdots \cup T_r]$
modulo $q$, and for each $i=1,\dots,r$ the same degree into $T_i$ modulo $q$.
\end{proposition}

\begin{proof}
Consider the map
\begin{align*}
  \phi \colon S &\to (\Z/q\Z)^{r+1}, \\
  \phi(v)&=\bigl(
    \deg_{G[S \cup T_1 \cup \cdots \cup T_r]}(v)\bmod q,\,
    \deg_{T_1}(v)\bmod q,\,
    \dots,\,
    \deg_{T_r}(v)\bmod q
  \bigr).
\end{align*}
Its codomain has exactly $q^{r+1}$ elements. Since $q^{r+1}k \le |S|$, the pigeonhole principle
yields a fiber $U \subseteq S$ with $|U| \ge k$ on which $\phi$ is constant. This is exactly the
stated residue constancy.
\end{proof}

Thus passing from one control block to $r$ control blocks costs only polynomially many residue
classes, namely $q^{r+1}$, and no additional difficulty in the transport step. The later frontier
theorem will show that, asymptotically, this extra bookkeeping does not create a new target.

\section{The canonical frontier}\label{sec:frontier}

We now define the two witness notions that matter most for the asymptotic story.

\begin{definition}
Let $G$ be a graph and $k \in \N$.
\begin{enumerate}
  \item A \emph{one-control exact witness of order $k$} consists of disjoint sets $S,T$ with
    $T \neq \varnothing$ and $|S| \ge k$ such that all vertices of $S$ have the same degree in
    $G[S \cup T]$ and the same degree into $T$.
  \item A \emph{one-control modular witness of order $k$} consists of disjoint sets $S,T$ with
    $T \neq \varnothing$, an integer $q \ge |S|$, and $|S| \ge k$ such that all vertices of $S$
    have the same degree in $G[S \cup T]$ modulo $q$ and the same degree into $T$ modulo $q$.
\end{enumerate}
\end{definition}

The exact one-control target looks stronger than the power-sequence target, but on geometric scales
it is already equivalent to it.

\begin{theorem}[single-control exact collapse]\label{thm:exact-collapse}
Fix an integer $b>1$. The following are equivalent.
\begin{enumerate}
  \item The power-sequence target holds at base $b$.
  \item For every $M \in \N$, every sufficiently large graph on $b^k$ vertices contains a
    one-control exact witness of order $Mk$.
\end{enumerate}

In fact, the control set may be chosen to have size $1$.
\end{theorem}

\begin{proof}
$(2)\Rightarrow(1)$ follows immediately from Proposition~\ref{prop:exact-transport}: every
one-control exact witness contains a regular induced subgraph on the surviving set $S$.

For $(1)\Rightarrow(2)$, assume the power-sequence target at base $b$. Fix $M \in \N$, and choose
$K$ such that every graph on $b^k$ vertices contains a regular induced subgraph on at least
$2Mk$ vertices for all $k \ge K$. Let $G$ be a graph on $b^k$ vertices with $k \ge K+1$, and
choose a vertex $t \in V(G)$.

Partition $V(G)\setminus\{t\}$ into the neighbors of $t$ and the non-neighbors of $t$. One of these
two sets has size at least $\lceil (b^k-1)/2 \rceil$. Since $b>1$ is an integer, this quantity is at
least $b^{k-1}$. Let $W$ be such a set, and choose a subset $W' \subseteq W$ of cardinality exactly
$b^{k-1}$.

By the power-sequence target applied to the induced subgraph $G[W']$, there exists $S \subseteq W'$
with $|S| \ge 2M(k-1)$ such that $G[S]$ is regular. Every vertex of $S$ has the same degree into the
single vertex $t$, namely either $0$ or $1$, because all vertices of $W$ lie in the same adjacency
class to $t$. Therefore $S$ together with the control set $\{t\}$ is a one-control exact witness of
order $2M(k-1)$ in $G$.

For $k \ge 2$ one has $2M(k-1)\ge Mk$, so this is in particular a one-control exact witness of order
$Mk$. This proves the required asymptotic statement.
\end{proof}

The modular route admits a larger family of witness notions. We will use the following four
variants.

\begin{definition}
Let $G$ be a graph and $k \in \N$.
\begin{enumerate}
  \item A \emph{one-control modular bucketing witness of order $k$} consists of pairwise disjoint
    sets $U,D,T$ with $T \neq \varnothing$, an integer $q \ge |U|$, and $|U| \ge k$ such that all
    vertices of $U$ have the same degree in $G[U \cup D \cup T]$ modulo $q$, the same degree into
    $D$ modulo $q$, and the same degree into $T$ modulo $q$.
  \item A \emph{nonempty modular control-block witness of order $k$} consists of a set $U$, a
    nonempty family of pairwise disjoint control blocks $T_1,\dots,T_r$ disjoint from $U$, an
    integer $q \ge |U|$, and $|U| \ge k$ such that all vertices of $U$ have the same degree in
    $G[U \cup T_1 \cup \cdots \cup T_r]$ modulo $q$, and for each $i$ they have the same degree into
    $T_i$ modulo $q$.
  \item A \emph{modular control-block bucketing witness of order $k$} consists of a set $U$, a
    dropped part $D$, and a nonempty family of pairwise disjoint control blocks
    $T_1,\dots,T_r$, all disjoint from one another, together with an integer $q \ge |U|$ and
    $|U| \ge k$, such that all vertices of $U$ have the same degree in
    $G[U \cup D \cup T_1 \cup \cdots \cup T_r]$ modulo $q$, the same degree into $D$ modulo $q$, and
    for each $i$ the same degree into $T_i$ modulo $q$.
  \item A \emph{modular control-block cascade witness of order $k$} consists of nested sets
    \[
      S_0 \supseteq S_1 \supseteq \cdots \supseteq S_m,
    \]
    a nonempty family of pairwise disjoint control blocks $T_1,\dots,T_r$ disjoint from $S_0$, and
    an integer $q \ge |S_m|$, such that $|S_m| \ge k$ and:
    \begin{enumerate}
      \item for each $i=1,\dots,m$, all vertices of $S_i$ have the same degree in
        $G[S_{i-1} \cup T_1 \cup \cdots \cup T_r]$ modulo $q$ and the same degree into the dropped
        layer $S_{i-1}\setminus S_i$ modulo $q$;
      \item all vertices of $S_m$ have the same degree in
        $G[S_m \cup T_1 \cup \cdots \cup T_r]$ modulo $q$, and for each $j$ they have the same
        degree into $T_j$ modulo $q$.
    \end{enumerate}
\end{enumerate}

For each of these witness notions, the corresponding \emph{target at base $b$} means that for every
$M \in \N$, every sufficiently large graph on $b^k$ vertices contains such a witness of order $Mk$.
\end{definition}

These variants differ only in whether one explicitly remembers the discarded vertices in a single
step or serializes them across a chain of refinements, and in whether the control data are stored in
one block or several. The next theorem shows that, asymptotically, they define the same frontier.

\begin{theorem}[canonical modular frontier]\label{thm:canonical-frontier}
Fix an integer $b>1$. The following are equivalent.
\begin{enumerate}
  \item The one-control modular target holds at base $b$: for every $M \in \N$, every sufficiently
    large graph on $b^k$ vertices contains a one-control modular witness of order $Mk$.
  \item The analogous one-control modular bucketing target holds.
  \item The analogous nonempty modular control-block target holds.
  \item The analogous modular control-block bucketing target holds.
  \item The analogous modular control-block cascade target holds.
\end{enumerate}
\end{theorem}

\begin{proof}
We prove the cycle
\begin{align*}
  &(1) \Rightarrow (2) \Rightarrow (1), \qquad
    (1) \Rightarrow (3) \Rightarrow (1), \\
  &(3) \Rightarrow (4) \Rightarrow (3), \qquad
    (4) \Rightarrow (5) \Rightarrow (4).
\end{align*}
Together these implications yield the equivalence of all five statements.

\smallskip

\noindent
$(1)\Rightarrow(2)$: every one-control modular witness is a special case of a one-control modular
bucketing witness, namely with empty dropped part $D=\varnothing$.

\smallskip

\noindent
$(2)\Rightarrow(1)$: let $U,D,T$ be a one-control modular bucketing witness. Since
\[
  \deg_{G[U \cup T]}(u)=\deg_{G[U \cup D \cup T]}(u)-\deg_D(u)
\]
for every $u \in U$, the ambient degrees in $G[U \cup T]$ are constant modulo $q$ on $U$ whenever
the ambient degrees in $G[U \cup D \cup T]$ and the degrees into $D$ are constant modulo $q$ on
$U$. The degrees into $T$ are already constant modulo $q$ on $U$, so $U,T,q$ form a one-control
modular witness.

\smallskip

\noindent
$(1)\Rightarrow(3)$: a one-control modular witness is a nonempty modular control-block witness with
exactly one control block.

\smallskip

\noindent
$(3)\Rightarrow(1)$: let $U,T_1,\dots,T_r$ be a nonempty modular control-block witness, and write
$T=T_1 \cup \cdots \cup T_r$. Because the blocks are pairwise disjoint,
\[
  \deg_T(u)=\deg_{T_1}(u)+\cdots+\deg_{T_r}(u)
\]
for every $u \in U$. Since each summand is constant modulo $q$ on $U$, the total degree into $T$ is
also constant modulo $q$ on $U$. The ambient degree in $G[U \cup T]$ is already constant modulo
$q$, so $U,T,q$ form a one-control modular witness.

\smallskip

\noindent
$(3)\Rightarrow(4)$: every nonempty modular control-block witness is a modular control-block
bucketing witness with empty dropped part.

\smallskip

\noindent
$(4)\Rightarrow(3)$: let $U,D,T_1,\dots,T_r$ be a modular control-block bucketing witness. For every
$u \in U$,
\[
  \deg_{G[U \cup T_1 \cup \cdots \cup T_r]}(u)
  =
  \deg_{G[U \cup D \cup T_1 \cup \cdots \cup T_r]}(u)-\deg_D(u).
\]
Thus the ambient degree on $U \cup T_1 \cup \cdots \cup T_r$ is constant modulo $q$ on $U$ whenever
the larger ambient degree and the degree into $D$ are constant modulo $q$ on $U$. The block-degree
hypotheses are unchanged, so $U,T_1,\dots,T_r,q$ form a nonempty modular control-block witness.

\smallskip

\noindent
$(4)\Rightarrow(5)$: let $U,D,T_1,\dots,T_r$ be a modular control-block bucketing witness. Define a
one-step chain by
\[
  S_0 := U \cup D, \qquad S_1 := U.
\]
The hypotheses already say that all vertices of $S_1$ have the same degree in
$G[S_0 \cup T_1 \cup \cdots \cup T_r]$ modulo $q$ and the same degree into the dropped layer
$S_0 \setminus S_1 = D$ modulo $q$. To verify the terminal condition, note that for every $u \in U$,
\[
  \deg_{G[U \cup T_1 \cup \cdots \cup T_r]}(u)
  =
  \deg_{G[U \cup D \cup T_1 \cup \cdots \cup T_r]}(u)-\deg_D(u).
\]
Hence the terminal ambient degree is constant modulo $q$ on $U$, and the block-degree hypotheses are
unchanged. Thus $S_0 \supseteq S_1$ together with $T_1,\dots,T_r$ and $q$ is a modular
control-block cascade witness.

\smallskip

\noindent
$(5)\Rightarrow(4)$: let
\[
  S_0 \supseteq S_1 \supseteq \cdots \supseteq S_m = U
\]
be a modular control-block cascade witness with control blocks $T_1,\dots,T_r$ and modulus $q$.
Set
\[
  D := S_0 \setminus U.
\]
If $m=0$, then $D=\varnothing$, and the terminal condition already gives a modular control-block
bucketing witness with empty dropped part.

Assume now that $m\ge 1$. Since $U \subseteq S_1$, the first cascade step implies that all vertices
of $U$ have the same degree in $G[S_0 \cup T_1 \cup \cdots \cup T_r]$ modulo $q$. The terminal
condition implies that all vertices of $U$ have the same degree in
$G[U \cup T_1 \cup \cdots \cup T_r]$ modulo $q$. Therefore, for every $u \in U$,
\[
  \deg_D(u)
  =
  \deg_{G[S_0 \cup T_1 \cup \cdots \cup T_r]}(u)
  -
  \deg_{G[U \cup T_1 \cup \cdots \cup T_r]}(u)
\]
is constant modulo $q$ on $U$. The terminal block-degree hypotheses are unchanged, and
$G[U \cup D \cup T_1 \cup \cdots \cup T_r]=G[S_0 \cup T_1 \cup \cdots \cup T_r]$. Thus
$U,D,T_1,\dots,T_r,q$ form a modular control-block bucketing witness.
\end{proof}

The previous two sections combine to show that the modular frontier still implies the original
conjecture.

\begin{corollary}\label{cor:frontier-implies-target}
Fix an integer base $b>1$. If every sufficiently large graph on $b^k$ vertices contains a
one-control modular witness of order $Mk$ for every $M$, then $F(n)/\log n \to \infty$.
\end{corollary}

\begin{proof}
By Proposition~\ref{prop:modular-transport}, every one-control modular witness yields a modular
witness on the surviving set, and Proposition~\ref{prop:modular-collapse} turns that modular
witness into a regular induced subgraph of the same order. Theorem~\ref{thm:modular-reduction}
finishes the argument.
\end{proof}

There are also bounded forms worth recording. The first shows that, asymptotically, any eventually
positive control budget is as good as a singleton budget on the modular side.

\begin{corollary}[eventual positive budget]\label{cor:positive-budget}
Fix an integer base $b>1$, and let $u \colon \N \to \N$ satisfy $u(k) \ge 1$ for all sufficiently
large $k$. Then the following are equivalent:
\begin{enumerate}
  \item For every $M \in \N$, every sufficiently large graph on $b^k$ vertices contains a
    one-control modular witness of order $Mk$ with a control set of size at most $u(k)$.
  \item $F(n)/\log n \to \infty$.
\end{enumerate}
\end{corollary}

\begin{proof}
If (1) holds, then Corollary~\ref{cor:frontier-implies-target} applies a fortiori.

Conversely, assume $F(n)/\log n \to \infty$. Choose $K_0$ such that $u(k) \ge 1$ for all
$k \ge K_0$. By Theorem~\ref{thm:exact-collapse}, for every $M \in \N$ there exists $K_1$ such that
every graph on $b^k$ vertices with $k \ge K_1$ contains a one-control exact witness
$(S,\{t\})$ of order $Mk$. For $k \ge \max\{K_0,K_1\}$, the singleton control set respects the
budget $u(k)$. Choosing any modulus $q \ge |S|$ turns $(S,\{t\})$ into a one-control modular
witness of the same order.
\end{proof}

\begin{corollary}[singleton budget]\label{cor:singleton-budget}
Fix an integer base $b>1$. The following are equivalent:
\begin{enumerate}
  \item For every $M \in \N$, every sufficiently large graph on $b^k$ vertices contains a
    one-control modular witness of order $Mk$ with a control set of size exactly $1$.
  \item $F(n)/\log n \to \infty$.
\end{enumerate}
\end{corollary}

\begin{proof}
This is the special case of Corollary~\ref{cor:positive-budget} with $u(k)=1$.
\end{proof}

\begin{corollary}[sublinear budgets collapse modular to exact]\label{cor:sublinear-budget}
Fix an integer base $b>1$, and let $u \colon \N \to \N$ satisfy $u(k)=o(k)$ and $u(k) \ge 1$ for all
sufficiently large $k$. Then the following are equivalent:
\begin{enumerate}
  \item For every $M \in \N$, every sufficiently large graph on $b^k$ vertices contains a
    one-control modular witness of order $Mk$ with a control set of size at most $u(k)$.
  \item For every $M \in \N$, every sufficiently large graph on $b^k$ vertices contains a
    one-control exact witness of order $Mk$ with a control set of size at most $u(k)$.
  \item $F(n)/\log n \to \infty$.
\end{enumerate}
\end{corollary}

\begin{proof}
Corollary~\ref{cor:positive-budget} gives the equivalence of (1) and (3). Also, (3) implies (2) by
Theorem~\ref{thm:exact-collapse}, because the singleton control set is eventually allowed by the
assumption $u(k) \ge 1$.

It remains to prove (1)$\Rightarrow$(2). Fix $M \in \N$. Because $u(k)=o(k)$, there exists $K_0$
such that $u(k) < Mk$ for all $k \ge K_0$. By (1), there exists $K_1$ such that every graph on
$b^k$ vertices with $k \ge K_1$ contains a one-control modular witness $(S,T,q)$ of order $Mk$
with $|T| \le u(k)$. For $k \ge \max\{K_0,K_1\}$, we therefore have
\[
  |T| \le u(k) < Mk \le |S| \le q.
\]
Since all degrees into $T$ are congruent modulo $q$ and lie in the interval $\{0,1,\dots,|T|\}$ of
width $<q$, they are all equal; write this common value as $e$. By
Proposition~\ref{prop:modular-transport}, the degrees in $G[S]$ are pairwise congruent modulo $q$,
and Proposition~\ref{prop:modular-collapse} then shows that $G[S]$ is regular because $q \ge |S|$.
Let $d$ be the common degree in $G[S]$. Then every vertex of $S$ has degree $d+e$ in $G[S \cup T]$,
so $(S,T)$ is a one-control exact witness of order $Mk$ with control set size at most $u(k)$.

Finally, (2)$\Rightarrow$(1) is immediate by choosing any modulus $q \ge |S|$.
\end{proof}

\begin{theorem}[summary reformulation package]\label{thm:summary-package}
Fix an integer base $b>1$. The following are equivalent.
\begin{enumerate}
  \item $F(n)/\log n \to \infty$.
  \item The power-sequence target holds at base $b$.
  \item The power-sequence modular target holds at base $b$.
  \item The power-sequence low-degree modular target holds at base $b$.
  \item For every $M \in \N$, every sufficiently large graph on $b^k$ vertices contains a
    one-control exact witness of order $Mk$.
  \item The one-control modular target holds at base $b$.
  \item The analogous one-control modular bucketing target holds.
  \item The analogous nonempty modular control-block target holds.
  \item The analogous modular control-block bucketing target holds.
  \item The analogous modular control-block cascade target holds.
  \item For every $M \in \N$, every sufficiently large graph on $b^k$ vertices contains a
    one-control modular witness of order $Mk$ with a control set of size exactly $1$.
\end{enumerate}
\end{theorem}

\begin{proof}
The equivalence of (1), (2), (3), and (4) is Theorems~\ref{thm:power-sequence}
and~\ref{thm:modular-reduction}. The equivalence of (2) and (5) is
Theorem~\ref{thm:exact-collapse}. Statements (6)--(10) are equivalent by
Theorem~\ref{thm:canonical-frontier}, and (6) implies (1) by
Corollary~\ref{cor:frontier-implies-target}. Finally, (1) and (11) are equivalent by
Corollary~\ref{cor:singleton-budget}, while (5) trivially implies (6) by choosing any modulus
$q \ge |S|$. Together these implications identify all eleven statements.
\end{proof}

\begin{remark}
Theorem~\ref{thm:summary-package} is the sense in which the paper is intended as a genuine
reduction rather than a vocabulary change. It shows that several witness frameworks that might
naturally be pursued as separate programs --- larger control families, explicit dropped parts, or
cascades of refinements --- do not define distinct asymptotic problems. Any proof in one of those
settings can be normalized to the one-control modular target without asymptotic loss. Accordingly,
the more elaborate witness packages can matter only if they retain additional quantitative
information beyond the bare existence statements considered here. The mathematical content of the
paper is therefore the removal of these false frontiers.
\end{remark}

\section{Concrete open problems}

The reformulation above leaves one canonical equivalent problem and one stronger exact route.

\begin{problem}[canonical modular frontier]\label{prob:canonical-frontier}
Fix an integer base $b>1$. Is it true that for every $M \in \N$, every sufficiently large graph on
$b^k$ vertices contains a one-control modular witness of order $Mk$?
\end{problem}

By Theorem~\ref{thm:summary-package}, Problem~\ref{prob:canonical-frontier} is equivalent to the
regular induced subgraph conjecture. It is the normalized frontier isolated by the paper: any
existence proof that proceeds through multiblock, bucketing, or cascade witnesses can be converted
back to this one-control target without asymptotic loss.

A stronger route would be to force exact control data with sublinear control size.

\begin{problem}[sublinear exact control]\label{prob:small-control}
Is there a function $u(m)=o(m)$ such that every sufficiently large graph contains disjoint sets
$S,T$ with
\[
  |S| \ge m, \qquad |T| \le u(m),
\]
all vertices of $S$ having the same degree in $G[S \cup T]$, and all vertices of $S$ having the
same degree into $T$?
\end{problem}

By Proposition~\ref{prop:exact-transport}, a positive answer to
Problem~\ref{prob:small-control} would immediately imply the regular induced subgraph conjecture.
This stronger problem is not claimed equivalent to the conjecture; it is included because it
isolates one natural exact-control route beyond the canonical modular frontier. The modular route
suggests an intermediate target: replace exact equalities by congruences modulo a modulus $q$
comparable to $|S|$, or even more sharply by congruences modulo a modulus larger than the induced
degree range, and then appeal to Proposition~\ref{prop:modular-collapse}. In either form, the
obstacle is no longer definitional. The canonical open step is
Problem~\ref{prob:canonical-frontier}; Problem~\ref{prob:small-control} asks whether one can do
better by forcing a more rigid witness.

\medskip

To summarize: the conjecture is equivalent to linear growth along any geometric subsequence, to the
corresponding modular target, to the sharper low-degree modular target, to the exact one-control
target with a singleton control vertex, and to each of the modular-control packages listed in
Theorem~\ref{thm:summary-package}. Problem~
\ref{prob:canonical-frontier} records the normalized equivalent frontier, while
Problem~\ref{prob:small-control} records a stronger exact route. The finite lemmas that support
this viewpoint replace full neighborhood data by ambient and external degree data, and replace
exponentially many profile classes by polynomially many residue classes. That is the structural
contribution of the present paper.

\begin{thebibliography}{99}

\bibitem{AlonKrivelevichSudakov}
N.~Alon, M.~Krivelevich, and B.~Sudakov,
\newblock Large nearly regular induced subgraphs,
\newblock {\em SIAM Journal on Discrete Mathematics} \textbf{22} (2008), no.~4, 1325--1337.

\bibitem{BermanRadcliffeScottWangWargo}
D.~M.~Berman, A.~J.~Radcliffe, A.~D.~Scott, H.~Wang, and L.~Wargo,
\newblock All trees contain a large induced subgraph having all degrees $1 \pmod k$,
\newblock {\em Discrete Mathematics} \textbf{175} (1997), 35--40.

\bibitem{CardosoKaminskiLozin}
D.~M.~Cardoso, M.~Kami\'nski, and V.~V.~Lozin,
\newblock Maximum $k$-regular induced subgraphs,
\newblock {\em Journal of Combinatorial Optimization} \textbf{14} (2007), 455--463.

\bibitem{Caro}
Y.~Caro,
\newblock On induced subgraphs with odd degrees,
\newblock {\em Discrete Mathematics} \textbf{132} (1994), 23--28.

\bibitem{FajtlowiczMcColganReidStaton}
S.~Fajtlowicz, T.~McColgan, T.~Reid, and W.~Staton,
\newblock Ramsey numbers for induced regular subgraphs,
\newblock {\em Ars Combinatoria} \textbf{39} (1995), 149--154.

\bibitem{KrivelevichSudakovWormald}
M.~Krivelevich, B.~Sudakov, and N.~Wormald,
\newblock Regular induced subgraphs of a random graph,
\newblock {\em Random Structures \& Algorithms} \textbf{38} (2011), no.~3, 235--250.

\bibitem{Scott}
A.~D.~Scott,
\newblock Large induced subgraphs with all degrees odd,
\newblock {\em Combinatorics, Probability and Computing} \textbf{1} (1992), 335--349.

\end{thebibliography}

\section*{Acknowledgments}

{\sloppy
The results of this paper were extracted from a companion Lean formalization, available at
\url{https://github.com/Arthur742Ramos/RegularInducedSubgraph}, which is intended as verification
support rather than as part of the mathematical contribution. The authors thank the Lean community
for the \texttt{mathlib} infrastructure underlying that formalization.\par}

\end{document}
